\subsection{Analytic realization of the floor thickenings}
\label{subsec:analytic-thickenings}

Retain the divisorial ideals \(\mathscr I_k\), the thickenings \(S_k\),
the map \(\phi:S\to W\) from Proposition~\ref{prop:conormal-layers},
and the sheaves \(\mathscr B_k\) in
\eqref{eq:thickening-ring-surjection}.  We construct their analytic
structure from the already analytic source thickenings.

\begin{lemma}[Analytic base thickenings]\label{lem:analytic-thickenings}
For every \(k\geq1\), put
\[
 W_k=(|W|,\mathscr B_k).
\]
Then \(W_k\) is a complex space with canonical reduction \(W\), and
the canonical sheaf evaluation defines a proper surjective holomorphic
map \(\phi_k:S_k\to W_k\) with
\[
 (\phi_k)_*\OO_{S_k}=\OO_{W_k}.
\]
Its reduction is \(\phi\).  The quotients
\(\mathscr B_{k+1}\twoheadrightarrow\mathscr B_k\) define closed
analytic embeddings \(W_k\hookrightarrow W_{k+1}\).  They commute
with the embeddings \(S_k\hookrightarrow S_{k+1}\) and the maps
\(\phi_k\).
\end{lemma}
\begin{proof}
\par\medskip\noindent\textit{Canonical rings and their successive kernels.}\par
All direct images defining \(\mathscr B_k\) first use the continuous
map \(|\phi|:|S|\to|W|\).  The equality
\(\mathscr B_1=\OO_W\) is the canonical equality for the holomorphic
map \(\phi\).  Apply derived direct image for abelian sheaves to the
conormal sequence.  Its derived image on an \(\OO_S\)-module is the
underlying sheaf of the analytic higher direct image.  To see this
formally, an injective \(\OO_S\)-module is flasque: for open sets
\(V\subset U\), the extension-by-zero modules satisfy
\(j_{V!}\OO_V\hookrightarrow j_{U!}\OO_U\).  Adjunction identifies
\(\operatorname{Hom}_{\OO_S}(j_{U!}\OO_U,J)\) with
\(\Gamma(U,J)\), so injectivity of \(J\) makes the restriction of
sections surjective.  Flasque abelian sheaves are acyclic for continuous
direct image.  An injective
module resolution therefore computes both the analytic module higher
images and their underlying abelian-sheaf higher images.  Thus
\eqref{eq:conormal-vanishing} gives exact sequences
\begin{equation}\label{eq:analytic-thickening-sequence}
 0\longrightarrow \mathscr F_k:=\phi_*\mathscr E_k
 \longrightarrow\mathscr B_{k+1}
 \longrightarrow\mathscr B_k\longrightarrow0\qquad(k\geq1).
\end{equation}
The displayed quotients are ring maps; exactness here concerns their
underlying abelian sheaves.  Proper direct-image coherence for the
already holomorphic \(\phi\)
\cite[Theorem~1.1, p.~15-01]{Grothendieck1961TechniquesVIII}
makes \(\mathscr F_k\) a coherent
\(\OO_W\)-module.

In \(\OO_{S_{k+1}}\), the ideal \(\mathscr E_k\) is square-zero and
is annihilated by \(\mathscr I_1\): use
\(\mathscr I_k^2\subset\mathscr I_{2k}\subset\mathscr I_{k+1}\)
and \(\mathscr I_1\mathscr I_k\subset\mathscr I_{k+1}\).
Consequently the action of \(\mathscr B_{k+1}\) on the ideal
\(\mathscr F_k\) factors through \(\mathscr B_1=\OO_W\).
This is the \(\OO_W\)-module structure on the square-zero ideal used
in the following construction.

The composites of \eqref{eq:analytic-thickening-sequence} are
epimorphisms \(\mathscr B_k\to\OO_W\).  Their kernels are
\[
 \mathscr N_k=|\phi|_*(\mathscr I_1/\mathscr I_k),
 \qquad \mathscr N_k^k=0,
\]
by left exactness and \(\mathscr I_1^k\subset\mathscr I_k\).
Since \(W\) is reduced, \(\mathscr N_k\) is exactly the nilradical.
Each stalk of \(\mathscr B_k\) is local: an element lifting a unit in
\(\OO_{W,w}\) has an inverse by a finite geometric series in the
nilpotent kernel, whereas an element mapping to the maximal ideal is
not a unit.  Its residue field is \(\C\).  Thus \(W_k\) is already
a locally ringed space, and it remains to construct local analytic
presentations for it.

\par\medskip\noindent\textit{A local analytic presentation from source functions.}\par
Fix \(k\) and \(w\in W\).  Choose a local closed analytic embedding
of a neighborhood of \(w\) into an open polydisc.  Lift the finitely
many germs of its coordinate restrictions through
\(\mathscr B_k\to\OO_W\), represent the lifts on a common
neighborhood, and shrink the polydisc.  This gives a closed embedding
\(i:U\hookrightarrow D\subset\C^d\) and sections
\(b_1,\ldots,b_d\in\mathscr B_k(U)\) reducing to the coordinate
restrictions.  Put \(Z=S_{k,U}\), the open subspace of \(S_k\) with
underlying set \(\phi^{-1}(U)\).  By definition,
\(\mathscr B_k(U)=H^0(Z,\OO_Z)\).

For a possibly nonreduced complex space, a tuple of holomorphic
functions defines a holomorphic map to affine space, functorially on
structure sheaves \cite[Theorem~1.1, p.~10-02]{Grothendieck1961TechniquesIII}.
Apply this theorem to the \(b_j\).  Their point values lie in \(D\),
so the map factors as
\[
 g:Z\longrightarrow D.
\]
The same tuple theorem identifies its actual analytic restriction to
the reduction with \(i\phi\), because the reduced coordinate functions
are those of \(i\phi\).  In particular \(|g|=i|\phi|\).  The
restriction \(\phi^{-1}(U)\to U\) is proper and \(i\) is closed, so
\(g\) is proper; nilpotents do not change properness of the underlying
continuous map.

The holomorphic map \(g\) makes
\(\mathscr A=g_*\OO_Z\) a genuine \(\OO_D\)-algebra.  Its underlying
module is coherent by the proper direct-image theorem
\cite[Theorem~1.1]{Grothendieck1961TechniquesVIII}.  Directly from the
underlying maps, as sheaves of \(\C\)-algebras,
\[
 \mathscr A=i_*(\mathscr B_k|_U).
\]
Reduction gives an epimorphism
\(\mathscr A\to i_*\OO_U\).  It is \(\OO_D\)-linear because the
actual analytic restriction \(g|_S=i\phi\) identifies the reduced
action of every ambient holomorphic function with its restriction to
\(U\).  The target is the coherent quotient of \(\OO_D\) by the
ideal of \(U\).  Hence its kernel \(\mathscr N\) is a coherent
\(\OO_D\)-module.  It is also a nilpotent ideal.  Shrink \(D\) once
more and choose sections \(n_1,\ldots,n_r\) generating
\(\mathscr N\) as a module.  They are actual sections of \(\OO_Z\)
whose reductions vanish.  Since \(\OO_D\to i_*\OO_U\) is an
epimorphism, subtraction of an ambient lift of a reduction gives the
stalkwise module equality
\begin{equation}\label{eq:analytic-module-generation}
 \mathscr A=\OO_D\cdot1+
              \sum_{\nu=1}^r\OO_D\cdot n_\nu.
\end{equation}
No splitting of the reduction map is used here.

Apply the tuple theorem again to the \(b_j\) and \(n_\nu\).  It gives
\[
 H=(g,n_1,\ldots,n_r):Z\longrightarrow P=D\times\C^r.
\]
If \(j:U\hookrightarrow P\) is \(u\mapsto(i(u),0)\), then
\(|H|=j|\phi|\), since the \(n_\nu\) are nilpotent.  Thus \(H\)
is proper.  The sheaf \(\mathscr A'=H_*\OO_Z\) is a coherent
\(\OO_P\)-module and the unit is an algebra map
\[
 \theta:\OO_P\longrightarrow\mathscr A'.
\]
As sheaves of rings \(\mathscr A'=j_*(\mathscr B_k|_U)\), so its
stalk at \(j(u)\) is \(\mathscr B_{k,u}=\mathscr A_{i(u)}\), and
its stalk outside \(j(U)\) is zero.  The stalk identification uses
the cofinality of restrictions of ambient neighborhoods among
neighborhoods in the closed subspace \(U\).

The map \(\theta\) is surjective on every stalk.  At \(j(u)\), any
element of the target has, by \eqref{eq:analytic-module-generation},
the form \(a_0+\sum a_\nu n_\nu\) with
\(a_\nu\in\OO_{D,i(u)}\).  The projection of \(H\) to \(D\) is
\(g\), and the extra coordinate \(t_\nu\) pulls back to \(n_\nu\).
The element is therefore the image of the convergent germ
\[
 a_0(z)+\sum_{\nu=1}^r a_\nu(z)t_\nu\in\OO_{P,j(u)}.
\]
At all other stalks the target is zero.  The kernel
\(\mathscr Q=\ker\theta\) is a coherent analytic ideal
\cite[pp.~9-07--9-08]{Grothendieck1961TechniquesII}.  It defines a closed
complex subspace \(Y=V(\mathscr Q)\subset P\) by
\cite[Definition~2.2 and adjacent construction, p.~9-10]{Grothendieck1961TechniquesII}.
Its support is exactly
\(j(U)\), since the target stalk there surjects onto the nonzero ring
\(\OO_{U,u}\).  Under the homeomorphism with \(U\), its structure
sheaf is \(\mathscr B_k|_U\), and its reduction is \(U\).

Every germ of \(\mathscr Q\) pulls back to zero by the definition of
\(\theta\).  The closed-subspace factorization criterion
\cite[pp.~9-04--9-05]{Grothendieck1961TechniquesII} therefore factors
\(H\) holomorphically through \(Y\).  Its sheaf map, under the
identified structure sheaf, is the canonical evaluation
\[
 |\phi|^{-1}(\mathscr B_k|_U)\longrightarrow\OO_Z.
\]
Indeed every germ of \(\mathscr B_k|_U\) lifts through the surjective
\(\theta\) to an ambient germ, whose pullback is precisely its value
as a section on the inverse image.  This proves the assertion for every
germ, not only the chosen coordinates.

\par\medskip\noindent\textit{Gluing and closed transition maps.}\par
The ringed space \((|W|,\mathscr B_k)\) was fixed before any of the
local choices.  The presentations just constructed make each of its
restrictions an analytic space by the local definition and adjacent
construction in \cite[Definitions~2.1--2.2 and adjacent prose]{Grothendieck1961TechniquesII}.
On overlaps the identifications are
the identity on the same sheaf of \(\C\)-algebras, and hence are
analytic isomorphisms satisfying the cocycle condition
\cite[Definition~2.4 and Proposition~2.5]{Grothendieck1961TechniquesII}.
This proves that \(W_k\) is a complex space.  It retains the Hausdorff
topology of \(W\), and its reduction is canonically \(W\).
Write \(\varepsilon_k:|\phi|^{-1}\mathscr B_k\to\OO_{S_k}\) for
the global sheaf evaluation.  It is locally the holomorphic factorization
above, so it defines \(\phi_k\).  Its reduction is \(\phi\), and its
underlying map is \(|\phi|\); thus it is proper and surjective.
The direct-image equality is the definition of its target structure
sheaf, with the canonical unit as the identification.

For the closed transition, use the construction with \(H\) at level
\(k+1\) and restrict its coordinate functions to the closed subspace
\(S_{k,U}\).  They give another proper holomorphic map to the same
\(P\), with underlying map \(j|\phi|\).  Its unit is the composite
\[
 \OO_P\twoheadrightarrow j_*(\mathscr B_{k+1}|_U)
          \twoheadrightarrow j_*(\mathscr B_k|_U).
\]
The second quotient stays surjective under the closed embedding \(j\),
as seen on its stalks.  The last sheaf with this ambient action is
coherent by proper direct image.  The two units therefore have nested
coherent ideal kernels, so their analytic subspaces give the required
closed inclusion.  These local inclusions glue because their ring maps
are the fixed quotients.  Naturality of the evaluation maps gives
\[
 \varepsilon_k\circ|\phi|^{-1}(\mathscr B_{k+1}\to\mathscr B_k)
  =(\OO_{S_{k+1}}\to\OO_{S_k})\circ\varepsilon_{k+1},
\]
which proves commutativity with the source inclusions.  No retraction
\(W_k\to W\), Cartesian square, or realization of the infinite tower
as one formal completion is asserted or needed.
\end{proof}

\subsection{Extension of a prime-floor contraction}
\label{subsec:prime-floor-extension}

\begin{proposition}[Extension from a prime floor]\label{prop:prime-floor-extension}
Under the hypotheses of Proposition~\ref{prop:conormal-layers}, there
are a normal compact analytic space \(Z\), a closed embedding
\(W\hookrightarrow Z\), and a proper bimeromorphic morphism
\(F:X\to Z\) such that \(F|_S\) is the given map \(\phi\) followed
by that embedding and
\begin{equation}\label{eq:extension-fibers}
 X\setminus S\simeq Z\setminus W,\qquad
 F^{-1}(|W|)=|S|,\qquad
 |F^{-1}(w)|=|\phi^{-1}(w)|\quad(w\in W).
\end{equation}
The natural map \(\OO_Z\to F_*\OO_X\) is an isomorphism, and the
fibers are connected.  For a sufficiently divisible global index
\(\ell\) of \(S\), the actual line \(\OO_X(-\ell S)\) is
\(F\)-ample; in particular \(F\) is projective.  More generally an
actual rational line on \(X\) whose restriction to \(S\) is
\(\phi\)-ample is \(F\)-ample.  The target is in Fujiki's class
\(\mathcal C\).
\end{proposition}
\begin{proof}
Choose a positive global index \(\ell\) such that \(\ell S\) is
Cartier and \(\OO_S(-\ell S)\) is \(\phi\)-ample.  Put
\[
 A=S_\ell=\ell S,\qquad A'=W_\ell,
 \qquad f=\phi_\ell:A\longrightarrow A'.
\]
Here \(A\) is the full effective Cartier divisor, including its
nilpotents; its local equation on normal \(X\) is a nonzerodivisor.
Lemma~\ref{lem:analytic-thickenings} makes \(f\) a proper surjective
holomorphic map of complex spaces and gives the natural equality
\(f_*\OO_A=\OO_{A'}\).

Cartier multiplication gives
\(\mathscr I_\ell\mathscr I_j=\mathscr I_{j+\ell}\).  Hence, for
every integer \(n>0\), on the full space \(A\) one has
\begin{equation}\label{eq:cartier-negative-layer}
 \OO_A(-nA)
 =\mathscr I_{n\ell}\otimes\OO_A
 =\mathscr I_{n\ell}/\mathscr I_{(n+1)\ell}.
\end{equation}
Its finite filtration by
\(\mathscr I_j/\mathscr I_{(n+1)\ell}\), for
\(n\ell\leq j\leq(n+1)\ell\), has successive quotients
\(\mathscr E_j\) for \(n\ell\leq j<(n+1)\ell\).  All indices are
at least one.  These filtration terms are \(\OO_A\)-modules, because
\(\mathscr I_\ell\mathscr I_j\subset\mathscr I_{(n+1)\ell}\).

Let \(a:S\hookrightarrow A\) and \(i:W\hookrightarrow A'\) be the
closed reduction embeddings.  The compatibility in
Lemma~\ref{lem:analytic-thickenings} gives \(fa=i\phi\).
Pushforward by a closed embedding is exact, so the composition identity
for derived direct images of abelian sheaves gives
\[
 R^q f_*(a_*\mathscr E_j)=i_*R^q\phi_*\mathscr E_j=0\quad(q>0).
\]
The long exact sequences of the finite filtration in
\eqref{eq:cartier-negative-layer} therefore give
\begin{equation}\label{eq:cartier-negative-vanishing}
 R^q f_*\OO_A(-nA)=0\qquad(q>0,\ n>0).
\end{equation}
This uses no vanishing for \(\mathscr E_0\) and keeps the full Cartier
thickening throughout.

The actual line \(\OO_A(-A)\) restricts to
\(\OO_S(-\ell S)\).  The underlying map of \(f\) is \(\phi\),
so the reductions of their fibers, as closed subspaces of \(X\),
coincide: a reduced closed analytic subspace is determined by its support.
Thus \(\OO_A(-A)\) is ample on the reduction of every \(f\)-fiber.
For a line on a compact nonreduced complex space, ampleness on the
reduction implies ampleness on the whole space.  Indeed the positive
metric from \cite[Corollary~1.12]{Fujino2026Fujiki} has local strictly
plurisubharmonic weights in common ambient embeddings, by its Lemma~2.4.
The same weights define a metric on the original line, since the
pointwise absolute values of transition units are unchanged by
nilpotents.  The converse positive-metric criterion gives ampleness on
the full space.  The proper fiberwise criterion
\cite[Definition~3.1 and Remark~3.2]{Fujino2026Fujiki} now makes
\(\OO_A(-A)\) \(f\)-ample.

We apply the analytic blowdown criterion stated in
\cite[Theorem~4.2, p.~14]{DasHacon2024OnMMP}, attributed there to
\cite[Theorem~2, p.~495]{Fujiki1975}.  For a reduced complex space,
an effective Cartier divisor \(A\) with its full possibly nonreduced
structure, and a proper surjective holomorphic map \(f:A\to A'\),
this criterion assumes \(f\)-ampleness of \(\OO_A(-A)\) and
\[
 R^1f_*\OO_A(-nA)=0\qquad\hbox{for every }n>0.
\]
It supplies a blowing down: a complex target \(Z\) containing \(A'\)
as an embedded subspace, a proper surjective holomorphic map
\(F:X\to Z\) whose restriction to \(A\) is \(f\), and an analytic
isomorphism of the complements.  In its universal form it also supplies
the natural equality of sheaves of rings
\begin{equation}\label{eq:universal-blowdown-ring}
 F_*\mathscr S_{X,A,f}=\OO_Z.
\end{equation}
Here \(\mathscr S_{X,A,f}\) consists along \(A\) of the germs in
\(\OO_X\) whose restrictions lie locally in
\(\operatorname{im}(f^{-1}\OO_{A'}\to\OO_A)\), and is \(\OO_X\)
away from \(A\).  This is a subsheaf of rings; no module coherence of
this particular sheaf is used.  The hypotheses of the stated criterion
follow from \eqref{eq:cartier-negative-vanishing} and the preceding
ampleness argument.  No flatness or reducedness of \(A\) or \(A'\)
is required.

The embedded \(A'\) is closed because it is compact.  Its reduction
embeds the original \(W\) as a closed subspace of \(Z\).  The
compatible reductions give \(F|_S=\phi\) followed by this embedding.
Removing a subspace depends only on its support, so the complement
isomorphism is \(X\setminus S\simeq Z\setminus W\).  It follows
that \(F^{-1}(|W|)=|S|\), and the prescribed restriction to \(A\)
gives the fiber-set equalities in \eqref{eq:extension-fibers}.
Outside \(W\) the fibers are single reduced points.  These assertions
concern underlying fiber sets over \(W\); they require no Cartesian
square of the closed thickenings.  The fibers are connected because
the floor fibers are connected.  The target is compact as the image of
compact \(X\).

We next prove normality; it is not an extra conclusion assumed from the
blowdown criterion.  For every open \(V\subset Z\), a section
\(s\in\OO_X(F^{-1}V)\) restricts to a section on the open space
\[
 A\cap F^{-1}V=f^{-1}(V\cap A').
\]
The equality \(f_*\OO_A=\OO_{A'}\) makes this restriction the pullback
of a unique section on \(V\cap A'\).  Thus every germ of \(s\) along
\(A\) satisfies the defining condition of \(\mathscr S_{X,A,f}\).
The reverse inclusion follows from \(\mathscr S_{X,A,f}\subset\OO_X\).
Consequently
\[
 F_*\mathscr S_{X,A,f}=F_*\OO_X,
 \qquad F_*\OO_X=\OO_Z
\]
as sheaves of rings, the second equality by
\eqref{eq:universal-blowdown-ring}.  This reasoning uses equality of
the displayed open subspaces, not of any closed fiber products.

The equality first makes \(Z\) reduced.  The complement is dense:
the inverse image of a nonempty open of \(Z\) is a nonempty open of
\(X\), and therefore meets \(X\setminus S\).  The complement
isomorphism thus makes \(F\) bimeromorphic.  Let
\(\nu:Z^\nu\to Z\) be the finite analytic normalization.  The
finite normalization exists by
\cite[Part~B, Section~4, Corollaries~2--3]{Houzel1961}.  The
dominant bimeromorphic map from normal \(X\) lifts to \(Z^\nu\).
Pullback along this lift gives inclusions, injective on the dense common
complement,
\[
 \OO_Z\subset\nu_*\OO_{Z^\nu}\subset F_*\OO_X=\OO_Z.
\]
The normalization is consequently an isomorphism and \(Z\) is normal.
The natural direct-image equality also gives connected fibers by analytic
Stein factorization.

Finally consider the actual global line \(\OO_X(-\ell S)\).  Its
restriction is ample on every reduced \(F\)-fiber: this is clear off
\(W\), and over \(W\) the reduction is the same embedded reduced
space as that of the corresponding \(\phi\)-fiber.  The nonreduced
upgrade and the proper fiberwise criterion used above therefore make
this one line \(F\)-ample.  The same proof works for any actual ambient
rational line with \(\phi\)-ample restriction after clearing its
global index.  In particular it works for the negative full adjoint
in Proposition~\ref{prop:conormal-layers}.  Thus \(F\) is projective.
A projective resolution with smooth source of the compact K\"ahler
\(X\) is again compact K\"ahler, and its composite with \(F\) is
bimeromorphic.  Hence \(Z\) lies in Fujiki's class \(\mathcal C\).
\end{proof}

The proposition supplies the ambient step for floor data satisfying its
stated hypotheses, including the two actual ample-line conditions.  In
the threefold argument below it extends divisor-to-curve contractions.  In
Section~\ref{sec:reduction} it extends a selected floor contraction of
the supported fourfold program.  When that floor contraction cuts out
the selected ray, the exact fiber description keeps the same contracted
curves, and the actual relatively ample lines make the ambient step
projective.  Its target K\"ahler class is then established in the
supported-program argument.
